% Chapter 16: Barriers, lookbacks and the Brownian bridge.

A barrier option is a European option that is cancelled, or activated, when the underlying
reaches a level; a lookback option pays a function of the highest or lowest price. Both depend
on the path through its extremes, and for geometric Brownian motion the joint law of the
endpoint and an extreme is known in closed form (the reflection principle). Reflecting a
density in the barrier, the \emph{method of images}, prices every continuously monitored
barrier and lookback; an infinite sequence of reflections handles two barriers.

Real contracts are monitored at dates. Discrete monitoring changes prices at order $h^{1/2}$ in
the spacing $h$ of the dates: for the course's monthly down-and-out call it raises the price by
$42\%$. A shift of the barrier by $e^{\pm0.5826\,\sigma\sqrt h}$ captures the leading term
(Broadie, Glasserman and Kou); exact discrete prices follow from a recursion over the dates. In
Monte Carlo the same $h^{1/2}$ is the bias of checking a continuous barrier on the time grid,
and the Brownian bridge removes it.

Throughout, $W$ is a Brownian motion under the risk-neutral measure $\Q$ with a filtration
$(\F_t)$, the parameters $r$, $y$ and $\sigma>0$ are constants, and
\begin{equation}\label{eq:bar-X}
  S_t=S_0e^{X_t},\qquad X_t=\nu t+\sigma W_t,\qquad \nu=r-y-\tfrac12\sigma^2,\qquad \lambda=\frac{2\nu}{\sigma^2}.
\end{equation}
We write $\overline X_T=\max_{t\le T}X_t$, $\underline X_T=\min_{t\le T}X_t$ (the maximum and
minimum exist because the paths are continuous), $\overline S_T=S_0e^{\overline X_T}$ and
$\underline S_T=S_0e^{\underline X_T}$, and $p_T(x)=(\sigma\sqrt T)^{-1}N'\big((x-\nu T)/(\sigma\sqrt T)\big)$ for
the density of $X_T$. Numbers refer to the course's test case $S_0=K=100$, $T=1$, $r=1\%$,
$y=2\%$, $\sigma=20\%$, for which $\nu=-0.03$ and $\lambda=-1.5$.

\section{The contracts}\label{sec:bar-contracts}

Fix a payoff $g(S_T)$, here a call $(S_T-K)^+$ or a put $(K-S_T)^+$, a barrier $B>0$ and a
\emph{monitoring set} $\mathcal T\subseteq[0,T]$: either $[0,T]$ (continuous monitoring) or
dates $t_1<\dots<t_m$, by default $t_j=jh$ with $h=T/m$; as in the course, time $0$ is not a
monitoring date. A \emph{down} barrier, normally below the spot, is \emph{hit} if $S_t\le B$ for
some $t\in\mathcal T$; an \emph{up} barrier, normally above it, is hit if $S_t\ge B$ for some
$t\in\mathcal T$. The \emph{out} option pays $g(S_T)$ if the barrier is not
hit, the \emph{in} option pays $g(S_T)$ if it is. Two directions, in or out, call or put: these
are the eight standard barrier options of \textcite{rubinstein-reiner-1991}, who priced them all.
A \emph{rebate} $R$ may be added: we pay it at $T$ to the holder of an out option that has been
knocked out, or of an in option that has never been knocked in (paying it at the hitting time is
\exref{ex:bar-one-touch}). Since $\1_{\{\text{hit}\}}+\1_{\{\text{not hit}\}}=1$ path by path,
\begin{equation}\label{eq:bar-parity}
  V^{\mathrm{in}}+V^{\mathrm{out}}=e^{-rT}\EQ[g(S_T)]+Re^{-rT}
\end{equation}
in any model, provided the two options have the same barrier, the same monitoring set and the
same rebate; with different monitoring sets the identity fails (\exref{ex:bar-degenerate}).

Some cases are trivial. If the spot is on the far side of the barrier at inception, a
continuously monitored out option is knocked out at once, and so it is when $S_0=B$
(\cref{thm:bar-reflection}(iii)); with dates $t_1>0$ it is not: with $S_0=B=100$ and monthly
dates the down-and-out call is worth $3.23135$ (\cref{sec:bar-discrete}). If
$K\ge B$, the up-and-out call pays nothing, and so does the down-and-out put if $K\le B$.

\emph{Lookbacks} use the extremes over $\mathcal T\cup\{0\}$: the floating-strike put pays
$\overline S_T-S_T$ and the call $S_T-\underline S_T$; the fixed-strike call pays $(\overline S_T-K)^+$
and the put $(K-\underline S_T)^+$.

\section{The reflection principle with drift}\label{sec:bar-reflection}

\begin{theorem}[Reflection principle with drift]\label{thm:bar-reflection}
Let $T>0$ and use the notation \eqref{eq:bar-X}.
\begin{enumerate}[label=(\roman*)]
  \item For every stopping time $\tau$ of the filtration generated by $W$, the process
    $W^\tau_t=2W_{t\wedge\tau}-W_t$, equal to $W$ up to $\tau$ and reflected about $W_\tau$ afterwards,
    is a Brownian motion.
  \item For a level $b\ne0$ and a Borel set $A$ on the same side of $b$ as $0$,
    \begin{equation}\label{eq:bar-image}
    \begin{aligned}
      \Q\big(\overline X_T\ge b,\ X_T\in A\big)&=e^{\lambda b}\!\int_Ap_T(x-2b)\,\dd x&&\text{if }b>0,\ A\subseteq(-\infty,b),\\
      \Q\big(\underline X_T\le b,\ X_T\in A\big)&=e^{\lambda b}\!\int_Ap_T(x-2b)\,\dd x&&\text{if }b<0,\ A\subseteq(b,\infty).
    \end{aligned}
    \end{equation}
  \item With $s=\sigma\sqrt T$: for $b\ge0$,
    $\Q(\overline X_T\ge b)=N\big(\frac{\nu T-b}{s}\big)+e^{\lambda b}N\big(\frac{-\nu T-b}{s}\big)$, and for
    $b\le0$, $\Q(\underline X_T\le b)=N\big(\frac{b-\nu T}{s}\big)+e^{\lambda b}N\big(\frac{b+\nu T}{s}\big)$.
    These functions of $b$ are continuous, so the extremes have no atoms on either side of $0$,
    and $\Q(\overline X_T>0)=\Q(\underline X_T<0)=1$.
  \item For $b<0$ the \emph{knock-out density}
    $q_T(x)=\Q(\underline X_T>b,\,X_T\in\dd x)/\dd x$ is $p_T(x)-e^{\lambda b}p_T(x-2b)$ for $x>b$ and $0$
    for $x\le b$; for $b>0$ the same with $\overline X_T<b$ and $x<b$. It is positive on the
    surviving side and vanishes at $x=b$.
\end{enumerate}
\end{theorem}

\begin{proof}
(i) On $[0,T']$, $W^\tau=W^{\tau\wedge T'}$, so we may assume $\tau$ bounded. The \emph{strong Markov property}
(\citealp[Theorem~2.16]{morters-peres}, for stopping times of the right-continuous filtration
generated by $W$, which include ours; \citealp[Section~2.6]{karatzas-shreve} for general
filtrations) says that $\beta_t=W_{\tau+t}-W_\tau$ is a Brownian motion independent of $\F_\tau$; it
needs continuous paths and a stopping time, not merely a random time. The stopped path $W_{\cdot\wedge\tau}$ and $\tau$ are
$\F_\tau$-measurable, and
\[
  W_t=W_{t\wedge\tau}+\beta_{(t-\tau)^+},\qquad W^\tau_t=W_{t\wedge\tau}-\beta_{(t-\tau)^+}.
\]
Since $-\beta$ is also a Brownian motion independent of $\F_\tau$, the triples
$(W_{\cdot\wedge\tau},\tau,\beta)$ and $(W_{\cdot\wedge\tau},\tau,-\beta)$ have the same law, and so
do $W$ and $W^\tau$, which are the same measurable function of them.

(ii) \emph{No drift, $\sigma=1$, $b>0$.} Let $\tau_b=\inf\{t:W_t=b\}$, a stopping time because
$\{\tau_b\le t\}=\{\max_{u\le t}W_u\ge b\}$ by continuity. The paths of $W$ and $W^{\tau_b}$
agree up to $\tau_b$, so each reaches $b$ by time $T$ exactly when $\tau_b\le T$, and on that
event $W^{\tau_b}_T=2b-W_T$. By (i), the pair (maximum on $[0,T]$, value at $T$) has the same
law for both processes (it is a measurable function of the path), hence for $A\subseteq(-\infty,b)$
\[
  \Q\big(\overline W_T\ge b,\,W_T\in A\big)=\Q\big(\tau_b\le T,\,2b-W_T\in A\big)=\Q(2b-W_T\in A),
\]
because $2b-W_T\in A$ forces $W_T>b$, hence $\tau_b\le T$. The right side is
$\int_A\varphi_T(x-2b)\,\dd x$ with $\varphi_T$ the $\mathcal N(0,T)$ density.

\emph{Drift.} Let $\mu=\nu/\sigma$, so $X/\sigma=W+\mu t$. For a bounded Borel functional $F$ of the
path on $[0,T]$, \cref{thm:is-girsanov}(i), applied to $G(w)=F(w)\,e^{\mu w_T-\mu^2T/2}$, gives
$\EQ[F(W+\mu\,\cdot)]=\EQ\big[F(W)\,e^{\mu W_T-\mu^2T/2}\big]$. With
$F=\1\{\max_{[0,T]}w\ge b/\sigma,\ w_T\in A/\sigma\}$ and the driftless case (extended from sets
to nonnegative functions of $W_T$),
\[
  \Q\big(\overline X_T\ge b,X_T\in A\big)=\int_{A/\sigma}e^{\mu w-\mu^2T/2}\,\varphi_T\Big(w-\frac{2b}\sigma\Big)\dd w
  =\int_{A/\sigma}e^{2\mu b/\sigma}\,\varphi_T\Big(w-\frac{2b}\sigma-\mu T\Big)\dd w,
\]
by completing the square. Substituting $x=\sigma w$ and using $2\mu b/\sigma=\lambda b$ gives
\eqref{eq:bar-image}. For $b<0$ apply this to $-X$, whose drift is $-\nu$ and whose $\lambda$
is $-\lambda$, at the level $-b$.

(iii) $\Q(\overline X_T\ge b)=\Q(X_T\ge b)+\Q(\overline X_T\ge b,X_T<b)$, and by (ii) the second term is
$e^{\lambda b}\Q(X_T+2b<b)=e^{\lambda b}\Q(X_T<-b)$. The formula is continuous in $b>0$, so
$\Q(\overline X_T=b)=0$, and letting $b\downarrow0$ gives $\Q(\overline X_T>0)=N(\nu T/s)+N(-\nu T/s)=1$.
Likewise for the minimum.

(iv) Subtract \eqref{eq:bar-image} from the law of $X_T$. At $x=b$ the exponents of $p_T(b)$ and
$e^{\lambda b}p_T(-b)$ differ by
$\lambda b-\big[(b+\nu T)^2-(b-\nu T)^2\big]/(2\sigma^2T)=\lambda b-2b\nu/\sigma^2=0$. Positivity on the
surviving side is \cref{prop:bar-bridge-crossing}(i): the ratio of the two terms is a probability
$e^{-2b(b-x)/(\sigma^2T)}<1$.
\end{proof}

The image $e^{\lambda b}p_T(x-2b)$ is the density of a mirror particle started at $2b$, weighted so
that it cancels the source on the barrier (\cref{fig:bar-images}, left); Girsanov's theorem
supplies the weight. Equivalently, $q_T$ solves the forward Kolmogorov equation
$\partial_Tq=\frac12\sigma^2q''-\nu q'$ on the surviving side with $q_T(b)=0$ (\exref{ex:bar-pde}), the
course's derivation; the probabilistic proof replaces the uniqueness theorem that route needs. By
(iii), $\Q(\underline S_T\le95)=0.82717$ for the test case.

\section{Closed forms by images}\label{sec:bar-images}

The \emph{surviving side} of a barrier $B$ is $\mathcal A=(B,\infty)$ for a down barrier and
$\mathcal A=(0,B)$ for an up barrier. For a payoff $g\ge0$ let
\begin{equation}\label{eq:bar-G}
  G(x)=e^{-rT}\EQ\big[g(xe^{X_T})\,\1_{\mathcal A}(xe^{X_T})\big],
\end{equation}
the Black--Scholes price of the European claim $g(S_T)\1_{\{S_T\in\mathcal A\}}$ when the spot is
$x$. For calls and puts, $G$ is a combination of vanilla prices $C(x,k)$, $P(x,k)$ and of
cash-or-nothing prices $D^c(x,k)$, $D^p(x,k)$ (paying $1$ if $S_T>k$, respectively $S_T<k$), all
with the parameters $T,r,y,\sigma$:
\begin{center}
\small
\begin{tabular}{@{}lll@{}}
\toprule
claim & $K\ge B$ & $K<B$\\
\midrule
call, down barrier & $C(x,K)$ & $C(x,B)+(B-K)D^c(x,B)$\\
call, up barrier & $0$ & $C(x,K)-C(x,B)-(B-K)D^c(x,B)$\\
put, down barrier & $P(x,K)-P(x,B)-(K-B)D^p(x,B)$ & $0$\\
put, up barrier & $P(x,B)+(K-B)D^p(x,B)$ & $P(x,K)$\\
\bottomrule
\end{tabular}
\end{center}
(For instance $(s-K)^+\1_{\{s>B\}}=(s-B)^++(B-K)\1_{\{s>B\}}$ when $K<B$.)

\begin{proposition}[Single barriers by images]\label{prop:bar-images}
Let $g\ge0$ be Borel with $\EQ[g(S_T)]<\infty$, monitor continuously, and in (i) and (ii) let
$S_0\in\mathcal A$.
\begin{enumerate}[label=(\roman*)]
  \item The out option paying $g(S_T)$ costs
    \begin{equation}\label{eq:bar-out}
      V^{\mathrm{out}}=G(S_0)-\Big(\frac B{S_0}\Big)^{\lambda}G\Big(\frac{B^2}{S_0}\Big),
      \qquad\lambda=\frac{2(r-y)}{\sigma^2}-1 .
    \end{equation}
  \item The in option costs $e^{-rT}\EQ[g(S_T)]-V^{\mathrm{out}}$. A rebate $R$ at $T$ adds
    $Re^{-rT}\pi$ to the out option and $Re^{-rT}(1-\pi)$ to the in option, where
    $\pi=\Q(\text{barrier hit})$ is given by \cref{thm:bar-reflection}(iii) at $b=\log(B/S_0)$.
  \item With the table above, (i) and (ii) price the eight standard barrier options. If $S_0\notin\mathcal A$
    (including $S_0=B$), the out option is worth $Re^{-rT}$ and the in option is the vanilla.
\end{enumerate}
\end{proposition}

\begin{proof}
(i) Take a down barrier, $b=\log(B/S_0)<0$. On survival $X_T>b$, so the payoff is
$g(S_T)\1_{\{X_T>b\}}-g(S_T)\1_{\{X_T>b\}}\1_{\{\underline X_T\le b\}}$. The first term has price
$G(S_0)$. For the second, \cref{thm:bar-reflection}(ii), extended from sets to nonnegative
functions, gives
\[
  \EQ\big[g(S_T)\1_{\{X_T>b,\ \underline X_T\le b\}}\big]=e^{\lambda b}\int_b^\infty g(S_0e^x)p_T(x-2b)\,\dd x
  =e^{\lambda b}\,\EQ\big[g(S_0e^{2b}e^{X_T})\1_{\{X_T+2b>b\}}\big],
\]
and $S_0e^{2b}=B^2/S_0$, $S_0e^{X_T+2b}>B$ if and only if $(B^2/S_0)e^{X_T}>B$, and
$e^{\lambda b}=(B/S_0)^\lambda$. So the discounted second term is $(B/S_0)^\lambda G(B^2/S_0)$. Up
barriers are identical with $X_T<b$. (ii) is \eqref{eq:bar-parity}, and the rebate is paid on an
event of probability $\pi$ or $1-\pi$. (iii) The table is checked case by case, as in the example
below it; the degenerate case is \cref{thm:bar-reflection}(iii).
\end{proof}

The image term is the same European claim on a fictitious stock started at the mirror point
$B^2/S_0$ (the reflection of $S_0$ in $B$ on the log scale), weighted by $(B/S_0)^\lambda$. For the
test case, the down-and-out call with $B=95$ costs $3.75412$ against $7.36429$ for the call, and
the down-and-in call $3.61017$; the up-and-out call with $B=115$ costs only $0.44998$, because its
payoff is largest just below the barrier, where survival is least likely. \Cref{fig:bar-images}
(right) compares the knock-out density with simulated survivors. With $K\ge B$ and $y=0$,
\eqref{eq:bar-out} for the down-and-out call reads $C(S_0,K)-(B/S_0)^{2r/\sigma^2-1}C(B^2/S_0,K)$, the
formula of \textcite{merton1973}.

\begin{figure}[ht]
  \centering
  \includegraphics[width=\textwidth]{ch16-images}
  \caption{Test case, $B=90$. Left: the density $p_T$ of $X_T$, the negative image and their sum on $x>b$,
    the knock-out density (\cref{thm:bar-reflection}(iv)). Right: histograms of $X_T$ over the surviving
    paths among $4\times10^5$ (our implementation, numpy \texttt{default\_rng}, seed $16$), with continuous
    monitoring (exact, by the minimum of the bridge from $0$ to $X_T$, \cref{prop:bar-bridge-crossing})
    and with $12$ monthly dates; dashed, the image formula for the knock-out density.}
  \label{fig:bar-images}
\end{figure}

\begin{remark}[What the images need]\label{rem:bar-term-structures}
The proof used constant coefficients twice: the reflected path must be a Brownian motion with
the same drift up to the weight $e^{\lambda b}$, and $p_T(x-2b)$ must be a translate of
$p_T$. With time-dependent $r$, $y$ or $\sigma$, a mirror source of constant weight no longer
vanishes on the barrier at every time, and the method fails in general. One exception: if
$r(t)-y(t)-\frac12\sigma(t)^2$ is proportional to $\sigma(t)^2$, the log price is a Brownian motion
with constant drift run on the variance clock $\int_0^t\sigma^2$ (\cref{rem:ps-time-change}), and the
formulas hold in that clock. Exponential barriers reduce to constant ones
(\exref{ex:bar-moving}). Otherwise one simulates (\cref{sec:bar-bridge}) or uses grids
(\cref{ch:20-grids}).
\end{remark}

\section{Double barriers}\label{sec:bar-double}

A double knock-out option pays $g(S_T)$ if $L<S_t<U$ for all $t\le T$, where $L<S_0<U$. Let
$\ell=\log(L/S_0)<0<u=\log(U/S_0)$ and $w=u-\ell=\log(U/L)$. One image cancels the source on one
barrier but not on the other; reflecting the image in the other barrier, and so on, produces an
infinite series.

\begin{proposition}[Double barriers]\label{prop:bar-double}
\begin{enumerate}[label=(\roman*)]
  \item For $x\in(\ell,u)$,
    \begin{equation}\label{eq:bar-double}
      \frac{\Q\big(\ell<X_t<u\ \forall t\le T,\ X_T\in\dd x\big)}{\dd x}
      =q(x)=\sum_{k\in\Z}e^{\lambda kw}\Big[p_T(x-2kw)-e^{\lambda u}\,p_T(x-2u-2kw)\Big].
    \end{equation}
  \item Let $q_n$ be the sum over $|k|\le n$. For $x\in(\ell,u)$ and $n\ge0$,
    \[
      |q(x)-q_n(x)|\le\frac{4\,e^{|\nu|w/\sigma^2}}{\sigma\sqrt{2\pi T}}\sum_{j\ge n}e^{-2j^2w^2/(\sigma^2T)},
    \]
    so the series converges faster than geometrically, at the rate $e^{-2n^2w^2/(\sigma^2T)}$.
  \item With $G$ the price \eqref{eq:bar-G} of $g(S_T)\1_{\{L<S_T<U\}}$, the double knock-out costs
    \[
      V=\sum_{k\in\Z}\Big(\frac UL\Big)^{\lambda k}\bigg[G\Big(S_0\Big(\frac UL\Big)^{2k}\Big)
      -\Big(\frac U{S_0}\Big)^{\lambda}G\Big(\frac{U^2}{S_0}\Big(\frac UL\Big)^{2k}\Big)\bigg],
    \]
    and truncating at $|k|\le n$ costs at most $e^{-rT}w\sup_{(L,U)}g$ times the bound in (ii).
\end{enumerate}
\end{proposition}

\begin{proof}
(i) Take first $\nu=0$, $\sigma=1$, so $X=W$, $a=\ell<0<u$, $w=u-a$. Let $\sigma_1$ be the first
hitting time of $u$, $\sigma_2$ the first hitting time of $a$ after $\sigma_1$, $\sigma_3$ that of
$u$ after $\sigma_2$, and so on; define $\rho_1,\rho_2,\dots$ in the same way starting with $a$.
A continuous path alternates finitely often between $a$ and $u$ on $[0,T]$; if it reaches $u$
first and then $k-1$ more alternating levels, exactly $k$ of the $\sigma_n$ and, as
$\rho_j=\sigma_{j+1}$, $k-1$ of the $\rho_n$ are at most $T$; so
\[
  \1\{\text{exit before }T\}=\sum_{n\ge1}(-1)^{n+1}\big(\1_{\{\sigma_n\le T\}}+\1_{\{\rho_n\le T\}}\big),
\]
with partial sums in $\{0,1,2\}$, which allows dominated convergence. Now fix $A\subseteq(a,u)$
and write $c_i=u$ for odd $i$ and $c_i=a$ for even $i$. Reflecting the path after $\sigma_n$ about
$c_n$ (\cref{thm:bar-reflection}(i); the reflected path has the same $\sigma_1,\dots,\sigma_n$) gives
$\Q(\sigma_n\le T,W_T\in A)=\Q(\sigma_n\le T,W_T\in A_1)$ with $A_1=2c_n-A$, a set beyond $c_n$. A
path that is at $c_{n-1}$ at time $\sigma_{n-1}$ and ends beyond $c_n$ must reach $c_n$ in between,
so the event $\{\sigma_n\le T\}$ may be replaced by $\{\sigma_{n-1}\le T\}$. Reflecting after
$\sigma_{n-1}$ about $c_{n-1}$ moves the set beyond $c_{n-1}$, and so on down to $\sigma_1$, after which
the set lies above $u$ and $\{\sigma_1\le T\}$ can be dropped. Two consecutive reflections, about $a$
and then $u$, translate by $2w$, so
$\Q(\sigma_{2k}\le T,W_T\in A)=\Q(W_T\in A+2kw)$ and $\Q(\sigma_{2k+1}\le T,W_T\in A)=\Q(W_T\in 2u+2kw-A)$,
and symmetrically for the $\rho$'s with $a$ in place of $u$ and $-w$ in place of $w$. Collecting densities, with $2a=2u-2w$,
\[
  q(x)=\sum_{k\in\Z}\big[\varphi_T(x-2kw)-\varphi_T(x-2u-2kw)\big],
\]
where the series of Gaussians converges absolutely, which justifies the interchanges. The drift
enters through the factor $e^{\mu x-\mu^2T/2}$ of the proof of \cref{thm:bar-reflection}(ii), which
again turns each translate $\varphi_T(x-c)$ into $e^{\mu c}\varphi_T(x-c-\mu T)$; rescaling by
$\sigma$ gives \eqref{eq:bar-double}.

(ii) In the driftless unit-variance coordinates, $|x|<w$ on $(a,u)$. For $|k|=j+1\ge1$ both
$|x-2kw|$ and $|x-2u-2kw|$ are at least $2jw$ (for instance, for $k=-j-1$,
$x-2u+2(j+1)w>a-2u+2(j+1)w=2jw+(w-u)+(a-u+w)$, and $w-u=-a>0$, $a-u+w=0$), so each of the four terms
with $|k|=j+1$ is at most $\varphi_T(2jw)$. The Girsanov factor is at most $e^{|\mu|w}$ on
$(a,u)$. Rescaling gives (ii). (iii) Integrate $e^{-rT}g(S_0e^x)$ against
\eqref{eq:bar-double} term by term as in \cref{prop:bar-images}: the translate by $2kw$ is the
spot $S_0(U/L)^{2k}$, the translate by $2u+2kw$ the spot $(U^2/S_0)(U/L)^{2k}$.
\end{proof}

\textcite{kunitomo-ikeda-1992} derived the series for exponentially curved barriers. For the test
case with $L=80$, $U=120$, the terms with $|k|\le2$ give the density to machine precision
($q-q_1$ is $3.6\times10^{-6}$ at most, the bound $2.9\times10^{-3}$), the probability of
staying in the corridor for a year is $0.38057$, and the double knock-out call costs $1.00574$.
The rate $e^{-2n^2w^2/(\sigma^2T)}$ deteriorates when $\sigma^2T$ is large compared with $w^2$, and
then an eigenfunction expansion of the same density is the better series (\exref{ex:bar-sine}).


\section{Lookback options}\label{sec:bar-lookback}

For a lookback the extreme itself is the payoff, and its moments follow from the law of
\cref{thm:bar-reflection}(iii). Let $s=\sigma\sqrt T$ and, for real $\kappa$, $c$ and $d$,
\[
  J(\kappa,c,d)=\int_c^\infty e^{\kappa z}N\Big(\frac{d-z}s\Big)\dd z
  =\frac{e^{\kappa d+\kappa^2s^2/2}N\big(\frac{d-c}s+\kappa s\big)-e^{\kappa c}N\big(\frac{d-c}s\big)}{\kappa}
\]
for $\kappa\ne0$ (integrate by parts), and $J(0,c,d)=s\,N'\big(\frac{d-c}s\big)+(d-c)N\big(\frac{d-c}s\big)$,
the limit $\kappa\to0$.

\begin{proposition}[Lookback options]\label{prop:bar-lookback}
\begin{enumerate}[label=(\roman*)]
  \item For $c\ge0$ and real $\kappa$,
    $\EQ\big[e^{\kappa\max(c,\overline X_T)}\big]=e^{\kappa c}+\kappa\big[J(\kappa,c,\nu T)+J(\kappa+\lambda,c,-\nu T)\big]$,
    and $\EQ\big[(e^{\overline X_T}-e^c)^+\big]=J(1,c,\nu T)+J(1+\lambda,c,-\nu T)$.
  \item The floating-strike lookback put, at inception, costs
    \[
      LP=S_0\Big[e^{-yT}N(d_+)\Big(1+\frac{\sigma^2}{2(r-y)}\Big)-e^{-yT}
      +e^{-rT}N(-d_-)\Big(1-\frac{\sigma^2}{2(r-y)}\Big)\Big],
    \]
    with $d_\pm=(r-y\pm\frac12\sigma^2)\sqrt T/\sigma$,
    for $r\ne y$, and the limit of this expression as $y\to r$ when $y=r$. When the maximum so far
    is $S^+\ge S_0$, the price is $S_0e^{-rT}\EQ[e^{\max(c,\overline X_T)}]-S_0e^{-yT}$ with $c=\log(S^+/S_0)$.
  \item The fixed-strike lookback call costs $S_0e^{-rT}\EQ[(e^{\overline X_T}-e^{c})^+]$ with
    $c=\log(K/S_0)$ if $K\ge S_0$, and $LP+S_0e^{-yT}-Ke^{-rT}$ if $K<S_0$.
  \end{enumerate}
The lookback call and the fixed-strike put follow from (i) applied to $-X$ (drift $-\nu$), since
$\underline X_T=-\max_{t\le T}(-X_t)$.
\end{proposition}

\begin{proof}
(i) For $f$ continuously differentiable and monotone,
\[
  f\big(\max(c,\overline X_T)\big)=f(c)+\int_c^\infty f'(z)\1_{\{\overline X_T>z\}}\,\dd z .
\]
Take expectations with Tonelli's theorem ($f'$ has one sign) and insert
$\Q(\overline X_T>z)=N\big(\frac{\nu T-z}s\big)+e^{\lambda z}N\big(\frac{-\nu T-z}s\big)$
(\cref{thm:bar-reflection}(iii)). With
$f(z)=e^{\kappa z}$ the two terms give $\kappa J(\kappa,c,\nu T)$ and $\kappa J(\kappa+\lambda,c,-\nu T)$; with
$f(z)=(e^z-e^c)^+$ on $z\ge c$ the same computation without the constant. The formula for $J$
is one integration by parts and the identity $\int_{\alpha}^\infty e^{\kappa sv}N'(v)\,\dd v=e^{\kappa^2s^2/2}N(\kappa s-\alpha)$.
(ii) The payoff is $S_0e^{\max(c,\overline X_T)}-S_T$, and $e^{-rT}\EQ[S_T]=S_0e^{-yT}$. Take $c=0$ and
$\kappa=1$. Since $\nu T+\frac12s^2=(r-y)T$, $\nu T/s=d_-$ and $1+\lambda=2(r-y)/\sigma^2$, we have
$(1+\lambda)\big(-\nu T+\frac12(1+\lambda)s^2\big)=(r-y)T$ and $-\nu T/s+(1+\lambda)s=d_+$, so
\[
  J(1,0,\nu T)=e^{(r-y)T}N(d_+)-N(d_-),\qquad
  J(1+\lambda,0,-\nu T)=\frac{\sigma^2}{2(r-y)}\big[e^{(r-y)T}N(d_+)-N(-d_-)\big],
\]
and collecting terms gives the formula. The price is continuous in $y$ by dominated convergence
($\overline X_T$ is continuous in $\nu$ and $e^{\overline X_T}\le e^{|\nu|T+\sigma\max_tW_t}$, which is
integrable by (i) with $\nu=0$), and $J(\kappa,c,d)$ is continuous at $\kappa=0$. (iii) If $K\ge S_0$,
the payoff is $S_0(e^{\overline X_T}-e^c)^+$; if $K<S_0$, it is $\overline S_T-K=(\overline S_T-S_T)+S_T-K$.
\end{proof}

For the test case the floating-strike put costs $17.23062$ and the call $14.27527$; the
fixed-strike call and put at $K=100$ cost $16.24550$ and $15.26039$, about twice the vanilla
prices. The floating-strike options were priced by \textcite{goldman-sosin-gatto-1979}, the
fixed-strike ones by \textcite{conze-viswanathan-1991}.

\section{Discrete monitoring}\label{sec:bar-discrete}

\subsection{Exact prices by recursion over the monitoring dates}

With dates $t_j=jh$, the knock-out indicator is a function of the vector
$(X_{t_1},\dots,X_{t_m})$, whose increments are independent $\mathcal N(\nu h,\sigma^2h)$
(\cref{prop:ps-gbm-exact}). The price is therefore an $m$-dimensional Gaussian integral; the course expresses it through $m$-variate normal distribution functions, which are expensive
beyond a few dates. A recursion integrates one date at a time.

\begin{proposition}[Discrete barriers by backward recursion]\label{prop:bar-discrete}
Let $I\subseteq\R$ be an open interval (the surviving set of $\log(S_t/S_0)$, for instance $(b,\infty)$,
$(-\infty,b)$ or $(\ell,u)$), and $g\ge0$ with $\EQ[g(S_T)]<\infty$. Define
$v_{m}(x)=g(S_0e^x)\1_I(x)$ and, for $j=m-1,\dots,1$,
$v_j(x)=\1_I(x)\int_Ip_h(x'-x)\,v_{j+1}(x')\,\dd x'$. The option paying $g(S_T)$ if
$X_{t_j}\in I$ for $j=1,\dots,m$ costs
\[
  V_m=e^{-rT}\int_Ip_h(x')\,v_1(x')\,\dd x' ,
\]
and, if $m\ge2$, $v_{m-1}(x)=\1_I(x)\,e^{rh}G_h(S_0e^x)$, where $G_h$ is the Black--Scholes price
\eqref{eq:bar-G} with maturity $h$ and surviving set $S_0e^I$.
\end{proposition}

\begin{proof}
By backward induction, $v_j(X_{t_j})=\EQ\big[g(S_T)\prod_{i\ge j}\1_I(X_{t_i})\,\big|\,\F_{t_j}\big]$:
for $j=m$ this is the definition, and given $\F_{t_j}$ the increment $X_{t_{j+1}}-X_{t_j}$ is
independent with density $p_h$, so the independence lemma and the tower property give the step.
The last line is $\EQ[v_1(X_{t_1})]$ discounted, and $v_{m-1}$ is a one-step European price.
\end{proof}

The payoff kink and the barrier enter only $v_{m-1}$, in closed form, so every integrand is
smooth on $I$, and composite Gauss--Legendre rules (panels of width $\sigma\sqrt h$, $8$ nodes each,
$I$ truncated at $8$ standard deviations $\sigma\sqrt T$) converge fast; one step is then a
matrix--vector product, the same for every date. For the course's monthly down-and-out call
($B=95$, $m=12$) this gives $5.329789841605$ in milliseconds with $232$ nodes, and doubling the
nodes leaves twelve digits unchanged. It reproduces all $112$ entries of the course's table of
discrete down-and-out calls ($B=85,\dots,100$, $m=1,2,3,4,6,12,\infty$) to the five decimals
shown, and bivariate normals check $m=2$ independently. \Cref{tab:bar-discrete} shows how slowly
the discrete prices approach the continuous one.

\begin{table}[ht]
\centering
\small
\caption{Discretely monitored down-and-out calls on the test case: exact prices
  (\cref{prop:bar-discrete}), continuity-corrected prices (\cref{prop:bar-bgk}) and the relative
  error of the latter; $m=\infty$ is \cref{prop:bar-images}. The last rows are the up-and-out call
  with $K=100$, $B=115$.}
\label{tab:bar-discrete}
\begin{tabular}{@{}llccccc@{}}
\toprule
& & $m=4$ & $m=12$ & $m=52$ & $m=252$ & $m=\infty$\\
\midrule
$B=95$ & exact & $6.22816$ & $5.32979$ & $4.57782$ & $4.14873$ & $3.75412$\\
& BGK & $6.06446$ & $5.30339$ & $4.57963$ & $4.14875$ & \\
& error & $-2.6\%$ & $-0.50\%$ & $+0.04\%$ & $+0.001\%$ & \\
$B=99$ & exact & $5.27574$ & $3.67565$ & $2.26089$ & $1.48813$ & $0.87818$\\
& BGK & $4.61914$ & $3.30925$ & $2.14258$ & $1.47501$ & \\
& error & $-12\%$ & $-10\%$ & $-5.2\%$ & $-0.9\%$ & \\

UOC & exact & $1.11608$ & $0.83710$ & $0.63269$ & $0.53166$ & $0.44998$\\
& BGK & $1.35690$ & $0.90999$ & $0.64809$ & $0.53468$ & \\
& error & $+22\%$ & $+8.7\%$ & $+2.4\%$ & $+0.57\%$ & \\
\bottomrule
\end{tabular}
\end{table}

\subsection{The continuity correction}

\begin{proposition}[Continuity correction; Broadie, Glasserman and Kou]\label{prop:bar-bgk}
Let $S$ be as in \eqref{eq:bar-X} with constant $r,y,\sigma>0$, let $h=T/m$, $t_j=jh$, and
$\beta=-\zeta(\frac12)/\sqrt{2\pi}=0.58260\ldots$, where $\zeta$ is the Riemann zeta function.
\begin{enumerate}[label=(\roman*)]
  \item \emph{Barriers.} Let $V(B)$ be the price of one of the eight standard barrier options
    with continuous monitoring and barrier $B\ne S_0$, and $V_m(B)$ the price of the same option
    monitored at $t_1,\dots,t_m$. As $m\to\infty$,
    \begin{align*}
      V_m(B)&=V\big(Be^{\beta\sigma\sqrt h}\big)+o(\sqrt h)&&\text{for an up barrier }(B>S_0),\\
      V_m(B)&=V\big(Be^{-\beta\sigma\sqrt h}\big)+o(\sqrt h)&&\text{for a down barrier }(B<S_0):
    \end{align*}
    the discrete contract is priced by the continuous formula with the barrier moved
    \emph{away from the spot} by the factor $e^{\beta\sigma\sqrt h}$.
  \item \emph{The size of the effect.} $V_m(B)-V(B)=\pm\beta\sigma\sqrt h\,BV'(B)+o(\sqrt h)$, with $+$
    for up and $-$ for down barriers. Whenever $V'(B)\neq0$, the continuous price, and the mean of
    any simulation that checks a continuous barrier only on a grid of step $h$, are off by an
    amount of exact order $h^{1/2}$.
  \item \emph{Extremes and lookbacks.} With $\widetilde M_m=\max_{0\le j\le m}X_{t_j}$,
    $\EQ\big[\widetilde M_m\big]=\EQ\big[\overline X_T\big]-\beta\sigma\sqrt h+O(h)$; and for $K>S_0$ the
    fixed-strike lookback call on $S_0e^{\widetilde M_m}$ costs
    $e^{-\beta\sigma\sqrt h}\,V^{\mathrm{lb}}(Ke^{\beta\sigma\sqrt h})+o(\sqrt h)$, where $V^{\mathrm{lb}}(k)$ is
    the continuous fixed-strike lookback call with strike $k$. So a maximum or a lookback
    inspected only on the grid also has a bias of exact order $h^{1/2}$.
\end{enumerate}
\end{proposition}

Part (i) is Theorem~1.1 of \textcite{broadie-glasserman-kou} for the down-and-in and
down-and-out calls and the up-and-in and up-and-out puts, and \textcite{kou-2003} extends it
to all eight options (with a simpler proof). Part (iii) is Lemma~3 and Theorem~4 of
\textcite{broadie-glasserman-kou-1999}; their Theorem~4 also shifts the prior maximum, which is
irrelevant here because $K>S_0$. Floating-strike lookbacks (their Theorem~3) are not part of
the proposition. BGK 1999, footnote~3, note that a continuous yield is accommodated by
$\nu=r-y$; the 1997 proof uses only that the drift is constant. Part (ii) follows
from (i) by a Taylor expansion, as $V$ is continuously differentiable in $B$ (differentiate \eqref{eq:bar-out}).
\Cref{X-ex:ps-richardson-sqrt} derives the expansion of $\EQ[\widetilde M_m]$ without drift from
Spitzer's identity.

\begin{remark}[Beyond geometric Brownian motion]\label{rem:bar-general}
\Cref{prop:bar-bgk} assumes exact grid values of a geometric Brownian motion. For a general
diffusion killed on leaving an open set and simulated by the Euler scheme with the killing
checked only at the grid dates, the weak error is also of order $h^{1/2}$, and this order is
exact; with the killing decided by the Brownian bridge between dates (the \emph{continuous} Euler
scheme) it is of order $h$, under conditions on the payoff \citep{gobet-killed-2000}.
\end{remark}

\emph{Why $\beta\sigma\sqrt h$ (heuristic, not a proof).} At the first date after a crossing the grid
value has overshot the barrier, by $\beta\sigma\sqrt h$ on average when the barrier is far away in
units of $\sigma\sqrt h$ (BGK 1997, Lemma~4.2): the discrete path sees a barrier that much further away.

\emph{When it fails.} The correction is asymptotic in $\log(S_0/B)/(\sigma\sqrt h)$: with $B=99$ the
barrier is $0.17$ monthly standard deviations from the spot, and the corrected price is $10\%$ too
low (\cref{tab:bar-discrete}; \textcite{broadie-glasserman-kou}, Section~2, note the same near the
spot). And the $o(\sqrt h)$ term can be large: the up-and-out call, whose payoff jumps from
$B-K=15$ to $0$ at the barrier, is overpriced by $8.7\%$ with monthly dates, the error decaying like
$h$ (\cref{fig:bar-discrete}). Unequal spacing, few dates and term structures are outside the
theorem; \cref{prop:bar-discrete} handles all three.

\begin{figure}[ht]
  \centering
  \includegraphics[width=\textwidth]{ch16-discrete}
  \caption{Test case. Left: down-and-out call against the barrier, exact for $m=4,12,52$ dates (solid),
    BGK (dashed; not at $B=S_0$) and continuous (grey). Right: absolute errors against $m$ of the
    continuous and BGK prices of the down-and-out call with $B=95$ (the BGK error changes sign between
    $m=16$ and $24$) and of the BGK price of the up-and-out call with $B=115$; reference slopes $-\frac12$
    and $-1$.}
  \label{fig:bar-discrete}
\end{figure}

\section{The Brownian bridge}\label{sec:bar-bridge}

Checking a continuous barrier only on the time grid of a simulation estimates the discrete price
$V_m(B)$, with a bias of order $h^{1/2}$ (\cref{prop:bar-bgk}(ii)): $+1.58$ with monthly steps and still
$+0.39$ with daily ones for the down-and-out call with $B=95$. The Brownian bridge
(\cref{prop:ps-bridge}) removes it: given the grid, the path between two dates is a bridge whose
extremes have a known law.

\begin{proposition}[Crossing probabilities of the bridge]\label{prop:bar-bridge-crossing}
Let $X$ be as in \eqref{eq:bar-X}, $0=t_0<t_1<\dots<t_m=T$ a grid (not necessarily uniform),
$h_j=t_j-t_{j-1}$, and $\mathcal G=\sigma(X_{t_0},\dots,X_{t_m})$.
\begin{enumerate}[label=(\roman*)]
  \item For $s\ge0$, $h>0$ and a level $b$,
    \[
      \Q\Big(\max_{s\le t\le s+h}X_t\ge b\ \Big|\ X_s=x,\ X_{s+h}=z\Big)=\exp\Big(-\frac{2(b-x)(b-z)}{\sigma^2h}\Big)
      \quad\text{if }x,z<b,
    \]
    and $\Q(\min_{[s,s+h]}X\le b\mid X_s=x,X_{s+h}=z)=\exp\big(-2(x-b)(z-b)/(\sigma^2h)\big)$ if $x,z>b$; the
    probability is $1$ if $x$ or $z$ is beyond $b$. It does not depend on the drift $\nu$. For the
    stock, the crossing probability of a barrier $B$ over a step from $S$ to $S'$ is
    $\exp\big(-2\log(B/S)\log(B/S')/(\sigma^2h)\big)$.
  \item Given $\mathcal G$, the events $\{\max_{[t_{j-1},t_j]}X\ge b\}$, $j=1,\dots,m$, are
    independent with the probabilities $p_j$ of (i) with $h=h_j$; in particular
    $\Q(\overline X_T<b\mid\mathcal G)=\prod_{j=1}^m(1-p_j)$, and likewise for minima.
  \item Given $\mathcal G$, the step maxima $M_j=\max_{[t_{j-1},t_j]}X$ are independent, and if
    $U_1,\dots,U_m$ are independent uniform variables on $(0,1]$, independent of $\mathcal G$, then
    \[
      \widehat M_j=\tfrac12\Big(X_{t_{j-1}}+X_{t_j}+\sqrt{(X_{t_j}-X_{t_{j-1}})^2-2\sigma^2h_j\log U_j}\Big)
    \]
    have the same joint law with $\mathcal G$ as the $M_j$ (for minima, take the minus sign before
    the square root). So $(X_{t_0},\dots,X_{t_m},\max_j\widehat M_j)$ has the law of
    $(X_{t_0},\dots,X_{t_m},\overline X_T)$.
  \item For a continuously monitored up-and-out claim with payoff $Y=e^{-rT}g(S_T)\1_{\{\overline X_T<b\}}$,
    $\EQ[g(S_T)^2]<\infty$, the estimator $\widehat Y=e^{-rT}g(S_T)\prod_j(1-p_j)$, computed from an exact
    sample of the grid values, is unbiased for $\EQ[Y]$ and $\Var\widehat Y\le\Var Y$; likewise for
    down barriers. Refining the grid
    (adding dates to it) can only increase $\Var\widehat Y$.
\end{enumerate}
\end{proposition}

\begin{proof}
(i) Restarting at $s$ (the process $X_{s+\cdot}-X_s$ is independent of $\F_s$ with the law of $X$)
reduces to $s=0$, $x=0$ and the level $b'=b-x>0$. For $z<b'$, \cref{thm:bar-reflection}(ii) gives
$\Q(\overline X_h\ge b',X_h\in\dd z)=e^{\lambda b'}p_h(z-2b')\,\dd z$, and dividing by $p_h(z)$,
\[
  \exp\Big(\lambda b'-\frac{(z-2b'-\nu h)^2-(z-\nu h)^2}{2\sigma^2h}\Big)
  =\exp\Big(\frac{2\nu b'}{\sigma^2}+\frac{2b'(z-\nu h)-2b'^2}{\sigma^2h}\Big)=\exp\Big(-\frac{2b'(b'-z)}{\sigma^2h}\Big),
\]
which is $(b-x)(b-(x+z))$ in the original variables. With $\pi(z)$ this expression for $z<b'$ and
$\pi(z)=1$ for $z\ge b'$, $\EQ[\1\{\overline X_h\ge b'\}\varphi(X_h)]=\EQ[\pi(X_h)\varphi(X_h)]$ for bounded Borel
$\varphi$, which is the statement. Minima follow by symmetry.

(ii)--(iii) For $t\in[t_{j-1},t_j]$ let $Y^j_t=X_t-\frac{(t_j-t)X_{t_{j-1}}+(t-t_{j-1})X_{t_j}}{h_j}$. The
drift cancels, $Y^j_t=\sigma\big(W_t-\text{its interpolation}\big)$, and by \cref{prop:ps-bridge}(i)
$Y^j_t$ is independent of $\mathcal G_{t_{j-1},t_j}=\sigma(W_r:r\notin(t_{j-1},t_j))$, which contains
$\mathcal G$ and every $Y^k_u$, $k\ne j$. All these variables are linear in $W$, hence jointly
Gaussian, and they are uncorrelated across the families $\mathcal G$, $Y^1,\dots,Y^m$; so the
processes $Y^1,\dots,Y^m$ and the grid vector are independent. Now $M_j=F(X_{t_{j-1}},X_{t_j},Y^j)$
for a measurable $F$ (a supremum over rational times, by continuity), and the independence lemma
gives the conditional independence, with conditional probabilities equal to versions of those
in (i). In (iii), $\widehat M_j$ solves $\exp(-2(\widehat M_j-X_{t_{j-1}})(\widehat M_j-X_{t_j})/(\sigma^2h_j))=U_j$ with
$\widehat M_j\ge\max(X_{t_{j-1}},X_{t_j})$, so by inversion (\cref{prop:rng-inversion}) its conditional law
is that of $M_j$ in (i).

(iv) By (ii), $\widehat Y=\EQ[Y\mid\mathcal G]$, and \cref{thm:st-conditional}(i) applies. On a coarser
grid contained in the finer one, the $\sigma$-algebra $\mathcal G$ is smaller, and
\cref{thm:st-conditional}(ii) gives the last claim.
\end{proof}

\begin{incode}
\coderef{code/src/mfe230d/barriers.py}: \texttt{barrier\_price}, \texttt{double\_barrier\_price},
\texttt{lookback\_floating}, \texttt{lookback\_fixed}, \texttt{discrete\_barrier\_price},
\texttt{bgk\_price}, \texttt{crossing\_probability}, \texttt{bridge\_extreme}, and
\texttt{barrier\_mc\_samples} with \texttt{monitor="grid"}, \texttt{"bridge"} or \texttt{"bridge-sample"}.
\end{incode}

\Cref{fig:bar-bridge} compares three estimators of the continuous down-and-out call ($B=95$, exact
$3.75412$). The bridge estimators are unbiased for every $m$, \emph{including $m=1$}: one normal for $S_T$ and the
weight $1-\exp(-2\log(S_0/B)\log(S_T/B)/(\sigma^2T))$ on $\{S_T>B\}$ give an unbiased estimate of the
continuously monitored price. The weighted estimator has the smallest variance with one step (standard
deviation $8.06$ per path, against $9.88$ with $12$ steps and $10.6$ for the sampled minima; our
implementation, numpy \texttt{default\_rng}, seed $7$, $2\times10^5$ paths), as (iv) predicts. Steps are needed when the payoff
looks at intermediate dates or the parameters change in time: \cref{prop:bar-bridge-crossing} holds
step by step when $r$, $y$, $\sigma$ are constant on each step, with $\sigma_j$ in place of $\sigma$ (same proof): the bridge survives the term structures that defeat the images.

\begin{figure}[ht]
  \centering
  \includegraphics[width=\textwidth]{ch16-bridge}
  \caption{Continuous down-and-out call, test case, $B=95$, from $m$ exact steps ($10^5$ paths per point;
    our implementation, numpy \texttt{default\_rng}, seeds $[16,i]$). Left: estimates $\pm2$ standard errors with grid monitoring
    and the two bridge estimators of \cref{prop:bar-bridge-crossing}; dashed, the closed form.
    Right: standard deviation per path.}
  \label{fig:bar-bridge}
\end{figure}

\emph{The course's discrete down-and-out call by simulation.} The monthly down-and-out call ($B=95$, price $5.32979$) needs no bridge, since it looks only at the
dates; but the continuous price, known in closed form, is a natural control variate.
\Cref{tab:bar-mc} reproduces the course's study, with coefficients estimated on the same paths
(\cref{thm:cv-optimal,thm:cv-multiple}). The continuous control is the sampled payoff
(\cref{prop:bar-bridge-crossing}(iii)), as in the course, or its conditional expectation given the
grid ((iv)), which is better correlated with the discrete payoff.

\begin{table}[ht]
\centering
\small
\caption{Monthly down-and-out call, $B=95$ (exact $5.32979$), $10^6$ paths or antithetic pairs (our
  implementation, numpy \texttt{default\_rng}, seed $2026$). VRF: variance reduction factor.}
\label{tab:bar-mc}
\begin{tabular}{@{}lcccc@{}}
\toprule
estimator & estimate & s.e. & VRF & course VRF\\
\midrule
plain & $5.33642$ & $0.01191$ & $1$ & $1$\\
antithetic pairs (\cref{prop:cv-antithetic}; $\rho=-0.201$) & $5.33197$ & $0.00752$ & $2.51$ & $2.51$\\
control: continuous DOC, sampled, $\beta^*=0.947$ & $5.33882$ & $0.00637$ & $3.49$ & $3.54$\\
control: continuous DOC, bridge weights, $\beta^*=1.101$ & $5.33621$ & $0.00491$ & $5.89$ & \\
control: vanilla call, $\beta^*=0.825$ & $5.32051$ & $0.00582$ & $4.19$ & $4.19$\\
both controls (sampled DOC, vanilla) & $5.32770$ & $0.00464$ & $6.59$ & $6.64$\\
both controls (bridge weights, vanilla) & $5.32900$ & $0.00399$ & $8.91$ & \\
\bottomrule
\end{tabular}
\end{table}

All estimates are within $1.6$ standard errors of the exact price. Antithetic pairs cost two
paths, so their efficiency gain (\cref{def:mc-efficiency}) is only about $1.25$.

\demo{16-barriers.ipynb}{images against simulated survivors, exact discrete prices against the
continuity correction, and grid against bridge Monte Carlo, for any parameters.}

\furtherreading
\textcite{glasserman-mc}, Section~6.4, treats barrier crossings by Brownian interpolation, and
\textcite{shreve2}, Chapter~7, barrier and lookback prices at the level of these notes.
